\section{Introduction}\label{sec:introduction}

Scientific information about a response often concerns its sign or how
it changes over an interval. Positivity restricts the response itself,
while monotonicity and convexity restrict its slope and curvature when
the corresponding derivatives exist. These properties need not be
visible in noisy observations, so a Gaussian process (GP) that describes
the data closely can still assign probability to curves that conflict
with the proposed scientific model. We therefore seek a posterior law
that imposes the restriction throughout a specified interval while
retaining joint uncertainty about the response and its derivatives.

The Hilbert space GP (HSGP) approximation of \citet{solin2020hsgp}
represents the process by a finite expansion in Laplacian eigenfunctions
on an interval extending beyond the region of interest. This extension,
called padding, is used to reduce boundary effects, while the rank,
or number of retained functions, controls the finest variation represented
for a given computational interval \citep{riutort2023practicalhsgp}.
The basis functions are independent of the covariance kernel and its
parameters. The kernel enters through spectral values that determine
Gaussian coefficient variances, allowing the same basis and its analytic
derivatives to be reused across covariance specifications.
Shape restrictions, however, concern
the entire curve obtained by summing the oscillatory basis functions.
Their validity is not determined by the signs of the individual
coefficients, and checking a finite set of locations can miss a
violation between them. The question addressed here is how to
condition this spectral model on the complete interval restriction
while retaining its basis and kernel-derived Gaussian coefficient law.

Our construction begins with positivity. A random intercept shifts
the whole curve above zero when it exceeds the negative minimum of
the spectral part. The same argument extends to monotonicity and
convexity whenever the retained span contains a direction with a
uniformly positive selected derivative, because a sufficiently large
coefficient along that direction makes the full curve satisfy the
restriction for any fixed values of the remaining coefficients. Thus the
continuum event becomes a lower bound on one Gaussian coefficient,
with a threshold determined by the remaining coefficients. Under
Gaussian observations, integrating out that coefficient gives an
explicit marginal posterior weighted by a normal survival probability.
A subsequent draw from its truncated conditional normal restores
the full conditioned law. This representation makes continuum
conditioning available within the HSGP expansion without
changing the probability model through a nonlinear transformation
of the curve.

\subsection{Related work and contributions}

Global derivative restrictions and derivative inference have precedents
in Bayesian regression, where piecewise-linear isotonic models permit
flat and increasing segments \citep{neelon2004isotonic}, while smooth
parameterizations and regression splines give other monotone priors
\citep{shively2009monotone}. Ordered Bernstein coefficients combine
global monotonicity with variable selection or clustering
\citep{curtis2011bernstein,wilson2020monotone}, with the aerosol analysis
of \citet{wilson2020monotone} illustrating derivative inference through
the derivative of a monotone differential-pressure curve.
For convex regression, maxima of random affine functions address
multivariate inputs and permit nonsmooth surfaces
\citep{hannah2011convex}, for which support and posterior consistency
results are available. Integrating squared GPs provides smooth
shape restrictions \citep{lenk2017shape}, with spectral constructions
also available for grouped data \citep{lenk2024hierarchical}.
These approaches establish that global shape information and
derivative inference can be incorporated in Bayesian regression in
several ways. They also make the choice of probability model
substantive: transforming or integrating a GP generally gives a
different law from conditioning a specified Gaussian model on a
shape event.

The closest constructions condition finite Gaussian expansions
using bases chosen to express global shape restrictions as coefficient
inequalities \citep{maatouk2017constraints}. These constraints hold
throughout the domain, and subsequent work develops general linear
inequalities, noisy observations, multiple constraints, and smooth
bases \citep{lopera2018finite,lopera2020noisy,maatouk2025bayesian}.
Within this line of work, \citet{ray2020shaperestricted} combine
smooth relaxations of coefficient constraints with elliptical slice
sampling and circulant embedding, while Gaussian simulation on
correlated subdomains addresses larger problems
\citep{maatouk2025largescale}. Our construction shares the aim of
conditioning a finite Gaussian model on a global restriction. Its
distinctive task is to characterize that restriction in the
Laplacian--Dirichlet spectral space, whose oscillatory basis requires
a different description of the constraint.

Other GP approaches impose derivative signs at virtual locations
\citep{riihimaki2010monotonicity,agrell2019operators}, use randomized
optimization for constrained inference \citep{zhang2026rlrto}, or
project GP curves onto the monotone-function space
\citep{lin2014projection}. Virtual-location restrictions provide a
finite description of the constraint, whereas projection defines
the fitted curve by a separate map. Here the target is the Gaussian
regression posterior conditioned on the complete continuum event
in the retained spectral space. The scalar threshold describes
that event exactly, so its numerical evaluation can be studied
separately from the choice of conditioning target.

The first contribution is a common representation for positivity,
monotonicity, and convexity, together with explicit directions and
sufficient rank and padding conditions for each case. Sign reversal
gives the corresponding nonpositive, nonincreasing, and concave
models. The representation also applies to joint restrictions whenever
the retained span contains a direction with the required signs at all
selected derivative orders. We give conditions on the HSGP direction
coefficients and threshold reductions covering all 26 combinations
without opposing constraints.

To evaluate the thresholds for the individual constraints, a cosine
transformation reduces the search for stationary
locations to polynomial equations, with endpoints, repeated roots,
and constant cases included. These characterizations make the
thresholds amenable to global evaluation. The resulting marginal
posterior is strongly log-concave in whitened coordinates, and a
total variation bound quantifies the change in its reconstructed
law when a conservative numerical threshold is used. The
computational cost is a global maximization at each target evaluation,
so the representation offers an alternative to bases with direct
coefficient inequalities, rather than a general efficiency advantage.

The second contribution connects finite continuum conditioning to
conditioning the underlying GP. Covariance approximation alone does
not settle this question because the indicator of a sign event is
discontinuous at curves that touch its boundary. For each individual
constraint, we establish a null boundary and control function-space
moments, obtaining quadratic Wasserstein convergence of the
conditioned posterior as padding and frequency cutoff grow. Under
the stated spectral regularity, convergence holds in the uniform
norm for the response and all derivatives up to the constrained
order. It therefore also gives uniform convergence of their posterior
means and covariance functions. The result includes Mat\'ern
smoothness greater than zero for positivity, greater than one for
monotonicity, and greater than two for convexity.

Two applications show why retaining joint uncertainty about the
response and its derivatives matters. Under the differential
absorption lidar approximation, mercury concentration is proportional
to the range derivative of a latent negative log signal ratio, so
nonnegative concentration requires that curve to be nondecreasing.
For leaf light-response data, concavity instead encodes a nonincreasing
marginal gain in assimilation as irradiance increases, without
requiring assimilation itself to increase throughout the range.
These analyses illustrate how a scientifically specified shape
restriction changes derivative inference.

This paper is organized as follows. \cref{sec:methodology} introduces the
spectral model and its conditioning events, followed by the Gaussian
coefficient construction for positivity, monotonicity, convexity, and
their joint combinations in \cref{sec:methods}. \cref{sec:theory} studies
posterior geometry, numerical conditioning error, and convergence to the
conditioned underlying GP, and \cref{sec:computation} describes posterior
computation. \cref{sec:applications} presents the mercury lidar and leaf
light-response applications, followed by the discussion in
\cref{sec:discussion}.
\ifdefined\CombinedInlineProofs\else
All proofs are given in the supplementary material.
\fi

Stan implementations of the monotone and convex models, together with
R code for reproducing the two applications, are available from the
author's GitHub repository.
